\documentclass[11pt]{article}
\usepackage[a4paper,margin=1in]{geometry}
\usepackage[T1]{fontenc}
\usepackage{lmodern,microtype}
\usepackage{amsmath,amssymb,amsthm,mathtools}
\usepackage{booktabs,enumitem}
\usepackage{xcolor}
\usepackage[numbers,sort&compress]{natbib}
\usepackage[colorlinks=true,linkcolor=blue!55!black,citecolor=blue!55!black]{hyperref}

\newtheorem{theorem}{Theorem}[section]
\newtheorem{corollary}[theorem]{Corollary}
\theoremstyle{remark}
\newtheorem{remark}[theorem]{Remark}

\title{A Uniform Trace-Envelope Criterion for Relative Regularized Determinants}
\author{Liang Wang}
\date{July 2026}

\begin{document}
\maketitle

\begin{abstract}
We formulate the all-order theorem missing from the finite trace atlas.  Let
$\tau_n(\sigma)$ be the cloud-extracted complementary power traces.  If
\[
 |\tau_n(\sigma)|\le Mq^n,\qquad n\ge2,
\]
uniformly in the noise, then for every $R$ with $Rq<1$ the logarithmic series
converges uniformly and
\[
 \sup_{|z|\le R}|\log R_\sigma(z)|
 \le M\bigl[-\log(1-Rq)-Rq\bigr]
 \le \frac{M(Rq)^2}{2(1-Rq)}.
\]
The relative regularized determinants are consequently locally bounded and
zero-free on that disk; Montel normality follows.

For the observed orders $2$--$12$, the smallest unit-amplitude rates read
directly from the maxima are $q_{\rm all}=0.35989$ and
$q_{\rm fine}=0.14377$.  The finite unit-disk log majorants are $0.07594$ and
$0.01074$.  If the same rates held at every later order, the all-order bounds
would be $0.10117$ and $0.01208$.  This last sentence is conditional: order
thirteen and above are not controlled by the audit.
\end{abstract}

\section{The trace formulation of the relative determinant}

For a Hilbert--Schmidt complement, the regularized determinant has germ
\[
 R_\sigma(z)=\exp\left(
 -\sum_{n=2}^{\infty}\frac{\tau_n(\sigma)}{n}z^n
 \right).
\]
The projection-free factorization of RH-234 identifies the same object from
the complementary eigenvalue product \citep{WangRH234,Simon2005}.  RH-235
shows why direct trace control can be much sharper than a Hilbert--Schmidt
norm bound \citep{WangRH235}.

\section{Uniform geometric envelope}

\begin{theorem}[Trace-envelope normality]\label{thm:envelope}
Let $\Sigma$ be an index set and let $\tau_n(\sigma)\in\mathbb C$ for
$\sigma\in\Sigma$ and $n\ge2$.  Suppose
\[
 |\tau_n(\sigma)|\le Mq^n
\]
for constants $M\ge0$ and $q\ge0$, uniformly in $\sigma$.  Then for every
$R\ge0$ with $Rq<1$, the series
\[
 L_\sigma(z)=-\sum_{n=2}^{\infty}\frac{\tau_n(\sigma)}{n}z^n
\]
converges uniformly on $|z|\le R$ and satisfies
\[
 |L_\sigma(z)|
 \le M\bigl[-\log(1-Rq)-Rq\bigr]
 \le \frac{M(Rq)^2}{2(1-Rq)}.
\]
Thus $R_\sigma(z)=e^{L_\sigma(z)}$ is a locally bounded zero-free normal
family on $|z|<q^{-1}$.
\end{theorem}

\begin{proof}
For $|z|\le R$,
\[
 \sum_{n=2}^{\infty}\frac{|\tau_n(\sigma)||z|^n}{n}
 \le M\sum_{n=2}^{\infty}\frac{(Rq)^n}{n}
 =M[-\log(1-Rq)-Rq]
 \le \frac{M}{2}\sum_{n=2}^{\infty}(Rq)^n.
\]
The geometric sum gives the displayed bound and the Weierstrass test gives
uniform convergence.  Exponentiation gives
$e^{-B}\le|R_\sigma(z)|\le e^B$, so the family and its reciprocals are
locally bounded.  Montel's theorem gives normality.
\end{proof}

\begin{corollary}[Locally uniform convergence criterion]
If, in addition, every fixed $\tau_n(\sigma)$ converges as
$\sigma\to0$, then $R_\sigma$ converges locally uniformly on
$|z|<q^{-1}$ to the determinant generated by the limiting traces.
\end{corollary}

\begin{proof}
Split the logarithmic series into a finite head and a uniformly small
geometric tail.  The finite head converges coefficientwise; then exponentiate.
\end{proof}

\begin{theorem}[Quantitative coefficient-to-determinant convergence]
Under the hypotheses of Theorem~\ref{thm:envelope}, suppose
$\tau_n(\sigma)\to\tau_n^{(0)}$ for each fixed $n$.  For $Rq<1$ and every
$N\ge2$,
\begin{align*}
 \sup_{|z|\le R}|L_\sigma(z)-L_0(z)|
 &\le \sum_{n=2}^{N}\frac{R^n}{n}
       |\tau_n(\sigma)-\tau_n^{(0)}|\\
 &\quad+
 \frac{2M(Rq)^{N+1}}{(N+1)(1-Rq)}.
\end{align*}
In particular the first line controls a finite coefficient head and the
second line is a uniform, explicitly summable tail.
\end{theorem}

\begin{proof}
The limiting coefficients inherit $|\tau_n^{(0)}|\le Mq^n$.  Split the
difference at order $N$ and bound each tail coefficient difference by
$2Mq^n$.  For $n\ge N+1$, use $1/n\le1/(N+1)$ and sum the geometric series.
\end{proof}

If $B_R=M[-\log(1-Rq)-Rq]$, the elementary exponential estimate also gives
\[
 \sup_{|z|\le R}|R_\sigma(z)-R_0(z)|
 \le e^{B_R}
      \sup_{|z|\le R}|L_\sigma(z)-L_0(z)|.
\]
Thus the theorem is suitable for a future validated coefficient audit: a
finite certified head plus an analytic tail bound produces a determinant
error bar on a compact disk.

\section{Observed finite envelope}

For a finite set of orders, define the unit-amplitude observed rate
\[
 q_{[2,m]}=\max_{2\le n\le m}
 \left(\sup_\sigma|\tau_n(\sigma)|\right)^{1/n}.
\]
This is a descriptive statistic, not an all-order certificate.

\begin{table}[ht]
\centering
\small
\begin{tabular}{lrr}
\toprule
Statistic & all endpoints & fine endpoints\\
\midrule
Observed orders & $2$--$12$ & $2$--$12$\\
$q_{[2,12]}$ & $0.359888$ & $0.143772$\\
Finite unit-disk majorant & $0.075936$ & $0.010737$\\
Conditional all-order bound & $0.101170$ & $0.012071$\\
First uncontrolled order & 13 & 13\\
\bottomrule
\end{tabular}
\caption{Finite observations and conditional geometric continuations.}
\end{table}

The finite majorants in the table sum the maximum at each order, even when
those maxima occur at different endpoints.  They are therefore slightly
larger than the largest single-endpoint jets of RH-236, namely $0.075920$ and
$0.010662$.  This distinction makes the table a valid uniform finite-order
majorant rather than a representative trajectory.

The global rate is set by the second moment.  Higher observed orders remain
subunit as well, but eleven points cannot exclude a later spike.  In
particular, fitting a line to the observed logarithms would not prove the
hypothesis of Theorem~\ref{thm:envelope}.

\section{What the criterion gives---and what it cannot identify}

The theorem controls both $R_\sigma$ and $R_\sigma^{-1}$.  Every locally
uniform subsequential limit on $|z|<q^{-1}$ is therefore holomorphic and
zero-free there.  This is exactly the analytic conclusion needed to prevent
the unresolved complement from creating or destroying a compact-set divisor
inside the controlled disk.

Normality is not identification.  Multiplying every member by a fixed
zero-free analytic factor whose logarithmic coefficients obey a compatible
geometric bound changes the limiting coefficients while preserving the same
type of criterion after a suitable adjustment of $M$.  Likewise a selector
that absorbs too many regular roots can produce an exceptionally small
relative factor.  The coefficient anchor must independently specify the
limiting logarithmic germ before the determinant can be matched to the
intended deterministic numerator.

The radius $q^{-1}$ is also local.  Nothing in the criterion continues the
limit across that boundary or supplies the global growth order needed for an
entire divisor comparison.  Those are later Gate-A questions even after a
unit-disk envelope is proved.

\section{Dynamical target}

The traces arise from noisy periodic-loop integrals at fixed positive noise.
The plausible analytic route is to extract the moving peripheral cloud in
the trace formula and prove a bound uniform simultaneously in orbit length
and noise.  Such a theorem must exploit dynamical cancellation; a generic
trace-ideal estimate returns to the RH-229 wall.

Concretely, one may seek a grouped periodic decomposition
\[
 \tau_n(\sigma)=\sum_{\mathcal G}W_{\mathcal G,n}(\sigma)
\]
in which the cloud contribution is removed before estimating the grouped
weights.  An absolute bound on every ungrouped loop may be too expensive and
recreate the Hilbert--Schmidt obstruction.  The useful theorem is instead a
cancellation-preserving majorant whose total is $O(q^n)$ uniformly as
$\sigma\to0$.

A proof may also be hybrid: certify the first $N$ traces directly, derive a
uniform long-orbit estimate for $n>N$, and use the quantitative theorem above
to join the two regimes.  This is more flexible than demanding a single
sharp formula at every order.

Even a successful envelope does not identify the limit by itself.  One also
needs a coefficient anchor proving that the selected cloud removes the
singular factor and leaves the intended deterministic numerator.  These are
the two next Gate-A walls.  Gates B--E and every Riemann-hypothesis conclusion
remain untouched.

\bibliographystyle{plainnat}
\bibliography{references}
\end{document}
